% preamble-exam-zh.tex — 共享 preamble
% 用于所有 exam-zh 模板
%
% 样式策略：
%   屏幕 → 语义彩色（蓝=关键想法，红=易错，绿=路线，灰=提示/教师备注）
%   打印 → 自动降级为黑白（通过 print mode 或 monochrome 选项）

\documentclass{exam-zh}

% ---------- 基础配置 ----------
\examsetup{
  page/size = a4paper,
  page/show-head = false,
  page/show-foot = false,
  sealline/show = false,
  font = times,
  math-font = xits,
}

\usepackage{tcolorbox}
\usepackage{needspace}
\tcbuselibrary{skins,breakable}

% ---------- 打印/屏幕颜色切换 ----------
% 用法：\documentclass[monochrome]{exam-zh} 即可切换为纯黑白
% 注意：不要使用 \if@xxx 放在 makeatother 外部；这里改成普通条件名。
\newif\ifeduprintcolor
\eduprintcolortrue

\makeatletter
\@ifclasswith{exam-zh}{monochrome}{\eduprintcolorfalse}{}
\makeatother

% ---------- 语义颜色定义 ----------
% 屏幕：有色彩区分；打印：降级为灰度
\ifeduprintcolor
  % --- 屏幕彩色模式 ---
  \definecolor{edu-blue}{RGB}{47, 82, 122}       % 关键想法
  \definecolor{edu-blue-bg}{RGB}{243, 246, 252}
  \definecolor{edu-red}{RGB}{180, 40, 40}         % 易错提醒
  \definecolor{edu-red-bg}{RGB}{255, 245, 240}
  \definecolor{edu-green}{RGB}{34, 100, 60}       % 路线图
  \definecolor{edu-green-bg}{RGB}{242, 250, 245}
  \definecolor{edu-orange}{RGB}{160, 90, 0}       % 训练目标
  \definecolor{edu-orange-bg}{RGB}{255, 248, 235}
  \definecolor{edu-gray}{RGB}{100, 100, 100}      % 提示/教师备注
  \definecolor{edu-gray-bg}{RGB}{248, 248, 248}
  \definecolor{edu-purple}{RGB}{90, 50, 120}      % 思考题
  \definecolor{edu-purple-bg}{RGB}{248, 244, 252}
\else
  % --- 打印黑白模式 ---
  \definecolor{edu-blue}{gray}{0.15}
  \definecolor{edu-blue-bg}{gray}{0.96}
  \definecolor{edu-red}{gray}{0.25}
  \definecolor{edu-red-bg}{gray}{0.97}
  \definecolor{edu-green}{gray}{0.20}
  \definecolor{edu-green-bg}{gray}{0.97}
  \definecolor{edu-orange}{gray}{0.25}
  \definecolor{edu-orange-bg}{gray}{0.98}
  \definecolor{edu-gray}{gray}{0.40}
  \definecolor{edu-gray-bg}{gray}{0.97}
  \definecolor{edu-purple}{gray}{0.30}
  \definecolor{edu-purple-bg}{gray}{0.98}
\fi

% ---------- 自定义教学环境 ----------
% 共性：enhanced + breakable（跨页不断裂）

% 原题卡片 — 强边框，突出显示
\newtcolorbox{problemcard}{
  enhanced, breakable,
  colback=edu-blue-bg,
  colframe=edu-blue,
  boxrule=1.2pt,
  arc=0pt,
  left=3mm, right=3mm, top=3mm, bottom=3mm,
}

% 解题路线图 — 左边线 + 浅底
\newtcolorbox{routebox}{
  enhanced, breakable,
  colback=edu-green-bg,
  colframe=edu-green!30,
  boxrule=0pt,
  borderline west={2.5pt}{0pt}{edu-green},
  left=3mm, right=3mm, top=3mm, bottom=3mm,
}

% 关键想法 — 蓝色左边线
\newtcolorbox{keyideabox}{
  enhanced, breakable,
  colback=edu-blue-bg,
  colframe=edu-blue!30,
  boxrule=0pt,
  borderline west={2.5pt}{0pt}{edu-blue},
  left=3mm, right=3mm, top=3mm, bottom=3mm,
}

% 易错提醒 — 红色左边线
\newtcolorbox{mistakebox}{
  enhanced, breakable,
  colback=edu-red-bg,
  colframe=edu-red!20,
  boxrule=0pt,
  borderline west={3pt}{0pt}{edu-red},
  left=3mm, right=3mm, top=2.5mm, bottom=2.5mm,
}

% 提示 — 灰色虚线左边线
\newtcolorbox{hintbox}{
  enhanced, breakable,
  colback=edu-gray-bg,
  colframe=edu-gray!20,
  boxrule=0pt,
  borderline west={2pt}{0pt}{edu-gray, dashed},
  left=3mm, right=3mm, top=2mm, bottom=2mm,
}

% 思考题 — 紫色虚线左边线
\newtcolorbox{thinkbox}{
  enhanced, breakable,
  colback=edu-purple-bg,
  colframe=edu-purple!20,
  boxrule=0pt,
  borderline west={2pt}{0pt}{edu-purple, dashed},
  left=2.5mm, right=2.5mm, top=2.5mm, bottom=2.5mm,
}

% 教师备注 — 灰色左边线，小字
\newtcolorbox{teachernote}{
  enhanced, breakable,
  colback=edu-gray-bg,
  colframe=edu-gray!15,
  boxrule=0pt,
  borderline west={2pt}{0pt}{edu-gray!60},
  left=2.5mm, right=2.5mm, top=2.5mm, bottom=2.5mm,
  fontupper=\small,
  colupper=black!60,
}

% 训练目标 — 橙色左边线，小字
\newtcolorbox{traininggoal}{
  enhanced, breakable,
  colback=edu-orange-bg,
  colframe=edu-orange!15,
  boxrule=0pt,
  borderline west={1.5pt}{0pt}{edu-orange},
  left=2mm, right=2mm, top=1mm, bottom=1mm,
  fontupper=\small,
}

% 答题区 — 无边框，纯留白
\newcommand{\answerarea}[1][40mm]{%
  \vspace{2mm}%
  \begin{tcolorbox}[
    enhanced, breakable,
    colback=white,
    colframe=white,
    boxrule=0pt,
    height=#1,
    valign=top,
    bottom=0mm,
    after skip=0mm,
  ]
  \end{tcolorbox}%
}

% ---------- 字体设置 ----------
% exam-zh (ctexbook) already sets CJK fonts via ctex.
% Only override if specific fonts are needed.
% macOS: Songti SC, STHeiti, STFangsong
% Linux: SimSun, SimHei, FangSong

% ---------- 页面设置 ----------
% exam-zh (ctexbook) already loads geometry.
% Override margins if needed:
% \geometry{top=16mm, bottom=16mm, left=14mm, right=14mm}

\examsetup{
  question/show-answer = true,
  fillin/show-answer = true,
  solution/show-solution = show-stay,
}

% ---------- 教师版解答样式 ----------
\newtcolorbox{solutionblock}{
  enhanced, breakable,
  colback=edu-blue-bg,
  colframe=edu-blue!25,
  boxrule=0pt,
  borderline west={2pt}{0pt}{edu-blue!50},
  arc=0pt,
  left=3mm, right=3mm, top=2mm, bottom=2mm,
  before skip=2mm,
  after skip=3mm,
  fontupper=\small,
}

% 解答步骤编号
\newcounter{solstep}
\newcommand{\solstep}[2]{%
  \stepcounter{solstep}%
  \par\medskip
  \noindent\hangindent=6mm
  {\color{edu-blue}\bfseries\textcircled{\scriptsize\thesolstep}\enspace #1}\enspace #2\par
}

\begin{document}

\section*{求一次函数解析式专题（三） · 练习卷（教师版）}

\section*{一、选择题（每题 12 分，共 48 分）}


\begin{keyideabox}
  \textbf{中等思想：对称直线解析式规律}
  已知直线 $y=kx+b$：\\ 1. 关于 $x$ 轴对称：用 $-y$ 替换 $y$ 得 $y = -kx - b$；\\ 2. 关于 $y$ 轴对称：用 $-x$ 替换 $x$ 得 $y = -kx + b$；\\ 3. 关于原点对称：两者都取反得 $y = kx - b$。
\end{keyideabox}



\begin{question}[points=12]
  直线 $y = -2x - 4$ 关于 $x$ 轴对称的直线解析式为
  \begin{choices}
    \item $y = 2x + 4$
    \item $y = -2x + 4$
    \item $y = 2x - 4$
    \item $y = -2x - 4$
  \end{choices}
  \paren[A]
  \begin{solutionblock}
    关于 $x$ 轴对称：$-y = -2x - 4 \implies y = 2x + 4$。
  \end{solutionblock}
\end{question}



\begin{question}[points=12]
  直线 $y = 3x - 6$ 关于 $y$ 轴对称的直线解析式为
  \begin{choices}
    \item $y = -3x + 6$
    \item $y = -3x - 6$
    \item $y = 3x + 6$
    \item $y = -3x$
  \end{choices}
  \paren[B]
  \begin{solutionblock}
    关于 $y$ 轴对称，用 $-x$ 替换 $x$ 得：$y = 3(-x) - 6 = -3x - 6$。
  \end{solutionblock}
\end{question}



\begin{question}[points=12]
  直线 $y = -x + 2$ 关于原点对称的直线解析式为
  \begin{choices}
    \item $y = -x - 2$
    \item $y = x - 2$
    \item $y = x + 2$
    \item $y = -x + 2$
  \end{choices}
  \paren[A]
  \begin{solutionblock}
    关于原点对称，横纵坐标皆取反：$-y = -(-x) + 2 \implies -y = x + 2 \implies y = -x - 2$。
  \end{solutionblock}
\end{question}



\begin{question}[points=12]
  已知一次函数 $y = kx + b$ 的图像与直线 $y = 5 - 4x$ 平行，并且它与直线 $y = -3(x - 6)$ 交于 $y$ 轴上的一点，则 $b$ 的值为
  \begin{choices}
    \item 18
    \item 5
    \item -18
    \item 6
  \end{choices}
  \paren[A]
  \begin{solutionblock}
    平行 $\implies k = -4$。交于 $y$ 轴上的点，即交点横坐标 $x = 0$。代入 $y = -3(x-6)$ 得 $y = 18$。所以交点为 $(0, 18)$。这说明 $y=kx+b$ 的截距 $b = 18$。
  \end{solutionblock}
\end{question}


\section*{二、填空题（每题 12 分，共 12 分）}


\begin{question}[points=12]
  直线 $y = 2x - 6$ 关于原点对称的直线解析式为 \fillin。
  \paren[$y = 2x + 6$]
  \begin{solutionblock}
    关于原点对称，用 $-y$ 替 $y$，$-x$ 替 $x$：$-y = 2(-x) - 6 \implies -y = -2x - 6 \implies y = 2x + 6$。
  \end{solutionblock}
\end{question}


\section*{三、解答题（共 40 分）}

\needspace{8\baselineskip}

\begin{problem}[points=20]
  将直线 $y = x - 2$ 向上平移若干个单位后，平移后的直线与反比例函数 $y = \dfrac{8}{x}$ 的图像在第一象限交于点 $C$。已知原直线与 $x$ 轴交于点 $A$，与 $y$ 轴交于点 $B$。若 $\triangle ABC$ 的面积为 $9$。

(1) 求点 $A$ 和点 $B$ 的坐标；
(2) 求平移后的直线解析式。
  \answerarea[60mm]
  \setcounter{solstep}{0}
  \begin{solutionblock}
    \textbf{答案：}(1) $A(2, 0)$，$B(0, -2)$；(2) $y = x + 7$\par\medskip
    \solstep{求 A、B 坐标，设解析式}{$A(2, 0)$，$B(0, -2)$。\\ 设向上平移 $a$ 个单位，解析式变为 $y = x - 2 + a$。\\ 令 $c = a - 2$，解析式为 $y = x + c$ ($c > -2$)。}
    \solstep{利用面积坐标公式列方程}{$C(x_c, y_c)$ 位于第一象限。$y_c = x_c + c$。\\ 根据三点坐标面积公式：$S_{\triangle ABC} = \frac{1}{2}|x_a(y_b - y_c) + x_b(y_c - y_a) + x_c(y_a - y_b)|$。\\ 代入 $A(2,0)$、$B(0,-2)$、$C(x_c, x_c+c)$：\\ $2S = |2(-2 - y_c) + 0 + x_c(0 - (-2))| = |-4 - 2(x_c+c) + 2x_c| = |-4 - 2c| = 2(c+2)$ (因为 $c > -2$)。\\ 故 $S_{\triangle ABC} = c + 2$。}
    \solstep{求解未知参数与解析式}{令 $S_{\triangle ABC} = 9 \implies c + 2 = 9 \implies c = 7$。\\ 故平移后的直线解析式为 $y = x + 7$。}
  \end{solutionblock}
\end{problem}


\needspace{8\baselineskip}

\begin{problem}[points=20]
  已知直线 $l: y = 2x - 4$。

(1) 分别写出直线 $l$ 关于 $x$ 轴对称的直线 $l_1$ 和关于 $y$ 轴对称的直线 $l_2$ 的解析式；
(2) 求直线 $l_1$ 与直线 $l_2$ 的交点坐标。
  \answerarea[50mm]
  \setcounter{solstep}{0}
  \begin{solutionblock}
    \textbf{答案：}(1) $l_1: y = -2x + 4$，$l_2: y = -2x - 4$；(2) 无交点，因为平行\par\medskip
    \solstep{写出对称线解析式}{关于 $x$ 轴对称的 $l_1$: 用 $-y$ 替 $y$ 得 $-y = 2x - 4 \implies y = -2x + 4$。\\ 关于 $y$ 轴对称的 $l_2$: 用 $-x$ 替 $x$ 得 $y = 2(-x) - 4 \implies y = -2x - 4$。}
    \solstep{求交点并分析}{联立 $y = -2x + 4$ 与 $y = -2x - 4$：\\ 得 $-2x + 4 = -2x - 4 \implies 4 = -4$ (方程无解)。\\ 所以这两条直线没有交点。\\ 原因是它们的斜率都为 $-2$，属于平行线关系。}
  \end{solutionblock}
\end{problem}


\clearpage
\section*{参考答案}


\begin{keyideabox}
  \textbf{选择题}
  1. A  2. B  3. A  4. A
\end{keyideabox}



\begin{keyideabox}
  \textbf{填空题}
  1. $y = 2x + 6$
\end{keyideabox}



\begin{keyideabox}
  \textbf{解答题}
  第 1 题：(1) $A(2, 0)$，$B(0, -2)$；(2) $y = x + 7$
第 2 题：(1) $l_1: y = -2x + 4$，$l_2: y = -2x - 4$；(2) 无交点，因为平行
\end{keyideabox}



\end{document}
